% =====================================================================================
% ARBITRAGE DES SUITES DES DEUX DERNIERS CALLS -- 13 aout 2026.
%
% PERIMETRE STRICT : uniquement les points souleves lors des deux derniers calls, celui avec
% Hugo et celui avec Caroline Hillairet. Le reste du memoire n'est PAS traite ici, meme quand
% il est concerne : ce document sert a trancher les suites d'un call, pas a relire le memoire.
%
% AJOUT DU 13 AOUT, SECONDE PASSE : l'allocation Euler/Shapley aux cinq piliers. Elle avait
% ete annoncee au call du 07 comme le sujet suivant et traitee en section C du point du 14,
% donc elle appartient bien au perimetre. Deux slides plutot qu'une, parce que ses DEUX
% RESERVES sont ce qui la rend utilisable : l'interdiction d'additionner les deux interactions
% (canaux et piliers), et l'instabilite de queue de l'allocation d'Euler a variance infinie.
% Les points suivants sont renumerotes de H5-H8 a H6-H9.
%
% CE QUE CHAQUE SLIDE DOIT PORTER, ET DANS CET ORDRE : ce qui a ete demande, ce que nous avons
% fait avec ses chiffres, puis une recommandation de GARDER ou non. La recommandation est le
% point neuf par rapport aux notes de suivi : celles-ci disaient ce qui etait fait, pas ce
% qu'il faut en conserver.
%
% AUCUNE FIGURE. Deux raisons : la regle du projet interdit de reprendre une figure deja
% servie dans un deck precedent, et un document de decision se lit mieux en tableaux qu'en
% graphiques. Cela evite aussi les debordements de cadre des figures portrait.
%
% DECOUPAGE EN \section* : sans cela le harnais lit un deck comme un bloc unique, et une
% declaration posee en tete avalerait tout le fichier. Les comptes de PAGES sont concentres
% sur une seule slide, elle-meme declaree, parce qu'ils viennent de main.toc et qu'aucun
% script ne les produit.
% =====================================================================================
\documentclass[11pt]{beamer}
\usetheme{Madrid}
\usecolortheme{seahorse}
\usepackage[utf8]{inputenc}
\usepackage[T1]{fontenc}
\usepackage[french]{babel}
\usepackage{booktabs}
\usepackage{amsmath,amssymb}
\usepackage{eurosym}
\setbeamertemplate{navigation symbols}{}
\setbeamertemplate{caption}{\raggedright\insertcaption\par}
\setbeamerfont{block title}{size=\small}

% SANS COULEUR, ET C'EST DELIBERE. Ces macros servent surtout en TITRE DE BLOC, or un titre de
% bloc beamer est deja du blanc sur un bandeau bleu fonce : un vert ou un orange y devient
% illisible. Le mot lui-meme porte l'information, et il reste lisible partout, en titre comme
% dans le corps du texte.
\newcommand{\garder}{\textbf{GARDER}}
\newcommand{\compresser}{\textbf{COMPRESSER}}
\newcommand{\annexer}{\textbf{DÉPLACER EN ANNEXE}}
\newcommand{\rien}{\textbf{RIEN À INTÉGRER}}
\newcommand{\atrancher}{\textbf{À TRANCHER PAR VOUS}}
\newcommand{\srcs}[1]{\vspace{0.8mm}{\tiny\textcolor{gray!50!black}{Sources : #1}}}

\title[Suites des deux derniers calls]{Suites des deux derniers calls}
\subtitle{Ce que nous avons fait sur chaque point, et ce qu'il faut en garder}
\author[K. Kaddouri]{Kélian Kaddouri}
\institute[ENSAE / Nexialog]{ENSAE / Nexialog Consulting}
\date[13/08/2026]{13 août 2026}

\begin{document}

\section*{Ouverture}

% HARNAIS-HORS-SECTION: page de titre et mode d'emploi. Millesimes et vocabulaire de
% recommandation, aucun resultat de modele.

\begin{frame}
\titlepage
\end{frame}

\begin{frame}{Ce que ce point traite, et ce qu'il ne traite pas}
\scriptsize
Ce document ne relit pas le mémoire. Il reprend \textbf{uniquement les points soulevés lors de
nos deux derniers calls}, un par slide, et pour chacun il dit trois choses dans cet ordre : ce
qui a été demandé, ce que nous avons fait avec ses chiffres, et \textbf{ce qu'il faut en
garder}.
\vspace{1.5mm}
\begin{block}{Le vocabulaire des recommandations}
\garder{} : reste tel quel. \quad \compresser{} : reste, mais réduit.\\
\annexer{} : le détail sort du corps, un paragraphe reste. \quad
\rien{} : la réponse existe, elle n'a pas à occuper de place.\\
\atrancher{} : le calcul est fait, la décision vous revient.
\end{block}
\vspace{-1mm}
\begin{block}{Pourquoi une recommandation, et pas seulement un état d'avancement}
Le corps du mémoire est aujourd'hui à \textbf{110 pages} contre les \textbf{70} que recommande
l'Institut. Répondre à une demande a donc un coût, et dire \og c'est fait \fg{} sans dire ce
que cela occupe serait incomplet. Une règle gouverne tous les arbitrages qui suivent : une
\textbf{limite déclarée} ne se retire jamais pour gagner une page. C'est ce qui tient ce
mémoire.
\end{block}
\end{frame}

\section*{Le coût en pages}

% HARNAIS-HORS-SCRIPT: comptes de pages lus dans main.toc apres compilation. Aucun script du
% pipeline ne les produit et aucun ne le devrait : ce sont des grandeurs de mise en page.

\begin{frame}{Ce que les deux calls ont produit, en pages}
\scriptsize
\begin{center}
\begin{tabular}{@{}l r l@{}}
\toprule
\textbf{Ce qui a été écrit} & \textbf{Pages} & \textbf{Où} \\
\midrule
Sensibilité aux probabilités de propagation & 2 & corps, résultats \\
Allocation aux cinq piliers, Shapley et Euler & 1 & corps, résultats \\
Défaillances simultanées et additivité & 2 & corps, robustesse \\
Séquences ordonnées, voie explorée et refermée & 2 & corps, identifiabilité \\
Les trois étages de l'incertitude & 2 & corps, robustesse \\
Scénarios par taille d'entité & 6 & corps, résultats \\
\midrule
Tableau des proxys & 2 & annexe \\
Martingalisation & 2 & annexe \\
Défaillances simultanées, le détail & 2 & annexe \\
Séquences ordonnées, le détail & 3 & annexe \\
Protocole d'élicitation & 3 & annexe \\
CTE agrégée & 1 & annexe \\
Niveau de l'intervalle reporté & 1 & annexe \\
\bottomrule
\end{tabular}
\end{center}
\vspace{1mm}
\begin{block}{Lecture}
\textbf{15 pages dans le corps, 14 en annexe.} Le corps est la ressource rare : c'est là que
porte l'arbitrage, et c'est pourquoi trois recommandations qui suivent proposent de déplacer
plutôt que de retirer. \textbf{Rien de ce qui suit ne propose de supprimer une limite
déclarée.}
\end{block}
\end{frame}

\section*{Les quatre demandes du call avec Hugo}

% HARNAIS-HORS-SCRIPT: la borne 0,84 est citee PRECISEMENT parce qu'aucun calcul du projet ne
% la reproduit. C'est le seul cas ou publier un nombre non confirme est correct : le fait qu'il
% soit introuvable EST l'information, et le harnais ne peut pas confirmer une absence.

\begin{frame}{H1. Retirer les conclusions tirées du bootstrap de $\xi$}
\scriptsize
\textbf{Demandé.} L'intervalle sur l'indice de queue est trop large pour porter une conclusion.
\vspace{1mm}

\textbf{Fait.} Les conclusions sont retirées. \textbf{L'intervalle reste affiché comme un fait}
et la limite déclarée ne bouge pas ; seule la lecture qu'on en faisait disparaît. Deux
trouvailles en faisant ce retrait :
\begin{itemize}\itemsep1pt
\item la slide portait $[0{,}32\,;0{,}84]$, un couple qu'\textbf{aucun calcul ne reproduit} : le
$0{,}32$ est la borne basse de la delta-méthode ($[0{,}320\,;0{,}871]$) et le $0{,}84$
n'appartient à aucun des deux intervalles. Le mémoire publie le bootstrap,
$[0{,}304\,;0{,}831]$ ;
\item le même travers vivait dans le mémoire, qui écrivait que $\hat\xi$ \og se stabilise
autour de $0{,}6$ \fg{} alors que le balayage le fait descendre à $0{,}38$.
\end{itemize}
\begin{block}{\garder}
Le retrait est acquis et ne coûte rien. Ce qui mérite d'être conservé au-delà de la demande :
le script imprimait en clair l'interdiction d'écrire cette phrase, et personne ne l'avait lue.
\end{block}
\srcs{script~\texttt{47}}
\end{frame}

\begin{frame}{H2. La table complète des leviers, et la réconciliation}
\scriptsize
\textbf{Demandé.} La somme des leviers pris un par un donnait $9\,138$~M\euro{} quand l'écart
total en vaut $14\,139$. D'où vient la différence.
\vspace{1mm}

\textbf{Fait.} Les seize configurations du treillis des quatre canaux sont énumérées. La
différence n'est pas une erreur, c'est l'\textbf{interaction}.
\begin{center}
\begin{tabular}{@{}l r r r@{}}
\toprule
\textbf{Canal} & \textbf{Isolé} & \textbf{Fermeture} & \textbf{Shapley} \\
\midrule
Fréquence & $4\,328$ & $8\,775$ & $6\,546$ \\
Détection & $1\,448$ & $4\,562$ & $2\,928$ \\
Accumulation & $1\,728$ & $3\,402$ & $2\,576$ \\
Propagation & $1\,633$ & $2\,402$ & $2\,088$ \\
\midrule
\textbf{somme} & $9\,138$ & $19\,141$ & $\mathbf{14\,139}$ \\
\bottomrule
\end{tabular}
\end{center}
\vspace{-1mm}
\begin{block}{\garder{} les trois colonnes}
L'interaction vaut $+5\,001$~M\euro, soit $+35\,\%$. \textbf{Les trois lectures ne sont pas
interchangeables et n'en publier qu'une détruirait le résultat} : un plan de remédiation doit
citer la colonne de \emph{fermeture}, parce qu'un canal fermé ferme aussi les croisés qu'il
portait.
\end{block}
\srcs{scripts~\texttt{43} et~\texttt{68}}
\end{frame}

\begin{frame}{H3. Les effets croisés entre leviers}
\scriptsize
\textbf{Demandé.} Regarder les croisements, en particulier fréquence et accumulation.
\vspace{1mm}

\textbf{Fait.} Décomposition de Möbius, quinze termes : ordre~1 $9\,138$, ordre~2 $5\,345$,
ordre~3 $-691$, ordre~4 $+346$, dont la somme redonne $14\,139$ à la précision machine. Les six
paires sont publiées avec leur bruit.
\vspace{1mm}

\textbf{Deux réserves, et elles comptent autant que le résultat.}
\begin{itemize}\itemsep1pt
\item \textbf{Il ne faut pas nommer de paire dominante} : les deux premières se tiennent à
$73$~M\euro{} pour des écarts-types de $708$ et $354$. Ce qui \emph{est} résolu : les trois
paires porteuses passent toutes par la fréquence et font $86\,\%$ de l'ordre~2 ;
\item la réconciliation est une \textbf{identité, pas une mesure}. Sur seize graines, neuf
termes sur quinze ont un signe résolu, et \textbf{aucun} des cinq d'ordre 3 et 4.
\end{itemize}
\begin{block}{\garder}
Le résultat à retenir est \emph{négatif} et c'est le plus solide : propagation $\times$
accumulation vaut $-330$~M\euro{} sur toutes les graines. Les deux canaux de co-occurrence sont
\textbf{substituts}, ce qui borne ce qu'on peut imputer à la contagion prise en général.
\end{block}
\srcs{script~\texttt{68}}
\end{frame}

\begin{frame}{H4. Le point 2 : portage, plancher, retour de 6 à 35 ans}
\scriptsize
\textbf{Demandé.} Les deux mots n'ont pas été compris, une clarification écrite est attendue.
\vspace{1mm}

\textbf{Fait, et la relecture a fait tomber une erreur.} Le mémoire écrivait que le
\emph{retour} était un \og plancher \fg{} alors que ce qui est minoré est le \emph{bénéfice}.
Or minorer le bénéfice \textbf{majore} la durée : le sens de la borne était inversé, dans le
mémoire, dans le script et dans le titre d'une figure. Les trois sont corrigés.
\begin{itemize}\itemsep1pt
\item \textbf{portage} : détenir du capital coûte chaque année. Libérer l'écart économise
$848$~M\euro/an à $6\,\%$, et $672$ au taux révisé de $4{,}75\,\%$ ;
\item \textbf{le bénéfice a deux composantes et une seule est monétisée} : la charge annuelle
passe de $3\,869$ à $622$~M\euro, soit une sinistralité évitée de $3\,247$~M\euro/an, qui vaut
$\mathbf{3{,}8}$ \textbf{fois} le portage ;
\item \textbf{d'où le sens} : le retour de $6$ à $35$ ans est un \textbf{MAJORANT}.
\end{itemize}
\begin{block}{\garder}
L'incompréhension en séance ne venait pas d'un défaut de pédagogie mais d'une erreur réelle. La
perte évitée, affirmée \og majoritaire \fg{} depuis le début sans jamais être calculée, est
désormais chiffrée.
\end{block}
\srcs{script~\texttt{50}}
\end{frame}

\section*{Les autres points du call avec Hugo}

\begin{frame}{H5. L'allocation aux cinq piliers : Shapley et Euler}
\scriptsize
\textbf{Annoncé au call comme le sujet suivant.} Attribuer le surcoût aux cinq piliers.
\vspace{1mm}

\textbf{Fait, et les deux méthodes répondent à deux questions différentes.} \emph{Qui porte le
surcoût} appelle le pilier \textbf{source}, et compose avec l'interaction : c'est Shapley.
\emph{Où loge le capital} appelle le pilier \textbf{touché}, à parts qui somment à cent : c'est
Euler.
\begin{center}
\begin{tabular}{@{}l r r r@{}}
\toprule
\textbf{Pilier} & \textbf{marginal} & $\boldsymbol{\varphi_j}$ \textbf{(Shapley)} & \textbf{part} \\
\midrule
P1 gouvernance & $3\,481$ & $\mathbf{3\,717}$ & $34\,\%$ \\
P4 tiers & $2\,658$ & $2\,810$ & $26\,\%$ \\
P2 incidents & $1\,634$ & $1\,867$ & $17\,\%$ \\
P3 tests & $1\,094$ & $1\,484$ & $14\,\%$ \\
P5 partage & $665$ & $1\,048$ & $10\,\%$ \\
\midrule
\textbf{somme} & $9\,531$ & $\mathbf{10\,927}$ & $100\,\%$ \\
\bottomrule
\end{tabular}
\end{center}
\vspace{-1mm}
\begin{block}{\garder}
\textbf{Le test non trivial n'est pas le classement mais la déformation.} Le classement
reproduit l'ordre des amorces, en partie mécaniquement. Ce qui ne l'est pas : la part de P1,
$34\,\%$, \textbf{excède} sa part d'amorce, $30\,\%$, et les $+1\,395$~M\euro{} redistribués
sont \textbf{intégralement} un produit de la cascade. Une page de corps pour une allocation
qu'un jury attend : le meilleur rapport du lot.
\end{block}
\srcs{scripts~\texttt{16b}, \texttt{20}, \texttt{20b}, \texttt{43} et~\texttt{67}}
\end{frame}

\begin{frame}{H5 \emph{(suite)}. Les deux réserves qui accompagnent l'allocation}
\scriptsize
\begin{block}{Ne jamais additionner les deux interactions}
Il y a celle entre les quatre \textbf{canaux}, $+5\,001$~M\euro{} de signe résolu, et celle
entre les cinq \textbf{piliers}, $+1\,395$~M\euro{} de signe \emph{non} résolu entre graines. Ce
sont \textbf{deux partitions du même écart}, pas deux effets qui s'ajoutent. Le titre d'une
slide du point précédent annonçait l'allocation comme \og la suite des $5\,001$ \fg, ce qui
invitait précisément à les additionner : il est corrigé.
\end{block}
\vspace{-1.5mm}
\begin{block}{Euler est stable en moyenne, instable en queue, et il faut le dire}
Les parts \emph{moyennes} sont stables, P2 $23\,\%$, P4 $22\,\%$, P1 $20\,\%$. Mais l'allocation
d'Euler repose sur des espérances \textbf{conditionnelles à la queue}, or à $\xi \approx 0{,}6$
la sévérité est de \textbf{variance infinie} : ces espérances convergent mal, et le classement
de queue bascule d'une résolution à l'autre. Je rapporte donc la queue d'un périmètre et tiens
l'autre pour indicative.
\end{block}
\vspace{-1.5mm}
\begin{block}{\garder{} les deux réserves}
\textbf{Ce sont elles qui rendent l'allocation utilisable.} Une allocation d'Euler publiée sans
mention de l'instabilité de queue serait plus fausse qu'utile, et deux interactions présentées
côte à côte sans avertissement s'additionnent dans la tête du lecteur. Aucune des deux ne coûte
plus de trois lignes.
\end{block}
\srcs{scripts~\texttt{20}, \texttt{20b} et~\texttt{68}}
\end{frame}

\begin{frame}{H6. Les scénarios par taille d'entité}
\scriptsize
\textbf{Demandé.} Chiffrer le coût de mise en conformité pour de petites, moyennes et grandes
entités.
\vspace{1mm}

\textbf{Fait, et le résultat est inconfortable.} Descendre le capital sectoriel au prorata de la
taille suppose une \textbf{élasticité de $1$}. Les deux canaux estimés en sont loin : fréquence
$+0{,}0744$, sévérité $0{,}087$ d'intervalle $[0{,}026\,;0{,}148]$, pour une élasticité globale
mesurée de $\mathbf{0{,}204}$. Une petite entité reçoit $0{,}633$ contre $0{,}100$ sous
proportionnalité, une grande $1{,}621$ contre $10{,}00$.
\vspace{1mm}

\textbf{Les deux approches se trompent en sens opposés} : la proportionnalité sous-estime une
petite entité, le modèle la majore parce qu'il ne plafonne pas la sévérité.
\begin{block}{\compresser}
Six pages du corps portent ce sujet. Le résultat utile tient en une page et demie : l'élasticité
mesurée, l'encadré des deux erreurs de sens opposés, et la requalification du chiffre d'entité
en \textbf{borne supérieure d'ordre de grandeur}. Le reste est de l'illustration.
\textbf{La limite déclarée, elle, reste} : c'est elle qui justifie la borne.
\end{block}
\srcs{scripts~\texttt{60} et~\texttt{70}}
\end{frame}

\begin{frame}{H7 et H8. Ne pas se brider, et la clarification écrite}
\scriptsize
\textbf{H7, demandé.} Ne pas se brider dans l'affichage de l'incertitude.
\vspace{1mm}

\textbf{Fait.} Toute grandeur simulée est désormais publiée \textbf{avec son bruit}, et les six
postures de report sont définies et chiffrées avec le leur : prédictive $\pm 7$,
$\beta = 90\,\%$ $\pm 13$, $\beta = 95\,\%$ $\pm 9$, $\beta = 99\,\%$ $\pm 24$. La table dit que
\textbf{la posture la plus prudente est la moins précise}, ce qui affaiblit en apparence la plus
rassurante.
\begin{block}{\garder}
Publier \og $1\,247 \pm 24$ \fg{} rend sans objet toute question sur le troisième chiffre.
Retirer le $\pm$ pour faire plus net serait exactement la dégradation contre laquelle ce mémoire
se construit. Réserve à tenir : sur quatre graines, l'ordre entre $\beta = 90$ et $95\,\%$
n'est \textbf{pas} résolu et n'a pas à l'être.
\end{block}
\vspace{-1mm}
\textbf{H8.} La clarification écrite sur le portage vous est due. Le contenu existe et est
vérifié, la slide H4 en est le résumé. \atrancher{} \; il ne reste qu'à l'envoyer.
\srcs{scripts~\texttt{46} et~\texttt{51}}
\end{frame}

\begin{frame}{H9. Quel écart entre états publier, et sous quelle définition}
\scriptsize
\textbf{Soulevé au call} en filigrane de l'analyse par levier, et non tranché depuis.
\vspace{1mm}

\textbf{Deux chiffres sont établis et imprimés, et ils ne mesurent pas la même chose.}
\begin{itemize}\itemsep1pt
\item un gain de conformité de $-53\,\%$, soit $8\,861$~M\euro{} de capital en moins, de
$16\,847$ à $7\,987$. Il fait varier le score de conformité \textbf{en gardant le gain de
propagation au niveau non conforme}, ce qui explique que son état \og conforme \fg{} soit à
$7\,987$ et non à $6\,049$ ;
\item un facteur $\mathbf{3{,}34}$ entre les deux états de référence, où \textbf{seuls les
quatre canaux bougent}, à sévérité de base et échelle inchangées.
\end{itemize}
\vspace{1mm}
Un troisième chiffre circule, de l'ordre de $-64\,\%$. \textbf{Il ne doit pas être employé en
l'état} : la vérification faite pour ce point ne le retrouve dans aucune sortie versionnée sur
cette grandeur.
\begin{block}{\atrancher}
\textbf{C'est la décision la plus urgente de ce document.} Tant qu'elle n'est pas prise, le
mémoire peut publier deux chiffres qui semblent se contredire alors qu'ils répondent à deux
questions différentes. Échanger l'un contre l'autre sans trancher remplacerait un chiffre mal
défini par un autre.
\end{block}
\srcs{scripts~\texttt{25} et~\texttt{43}}
\end{frame}

\section*{Le call avec Caroline : ce qui ne touche pas le mémoire}

\begin{frame}{C1 et C2. La dynamique de $Y_{j,t}$, et le mot \og gravité \fg}
\scriptsize
\textbf{C1, demandé.} Justifier la dynamique de $Y_{j,t}$, brownien pour le bruit de fond et
sauts pour les grands chocs.
\vspace{1mm}

\textbf{Fait.} Une note autonome de six pages. Trois horloges distinguées, et les deux échelles
sont \emph{identifiées} par la surdispersion mesurée : $1{,}19$ sur les cellules firme-année,
$2{,}24$ sur la série annuelle détendancée, contre $1$ pour un brownien pur.
\vspace{1mm}

\textbf{C2, demandé.} Clarifier l'emploi du mot \og gravité \fg.
\vspace{1mm}

\textbf{Fait, et le balayage donne l'inverse de ce qu'on attendait.} \textbf{Le mémoire
n'emploie nulle part \og gravité \fg{} pour le caractère systémique} : la confusion venait des
decks de juillet, archivés. En revanche il y a une \textbf{vraie collision dans le code}, entre
le gain de propagation et l'échelon de sévérité, dont les noms ne diffèrent que par un tiret
bas.
\begin{block}{\rien{} dans le mémoire}
Les deux réponses existent et ne demandent aucune page. La note sur $Y_{j,t}$ reste une pièce de
discussion, pas un chapitre. La collision de noms est un sujet de code, verrouillé en attendant
par la table des notations. Le renommage touche 27 fichiers sur un pipeline gelé : \atrancher.
\end{block}
\srcs{script~\texttt{08b}}
\end{frame}

\section*{Le call avec Caroline : les mesures demandées}

\begin{frame}{C3. Le test de martingalité}
\scriptsize
\textbf{Demandé.} Un test de martingalité est-il conduit, et sur quoi.
\vspace{1mm}

\textbf{Fait, et la réponse est un fait vérifié plutôt qu'un argument.} \textbf{Le mémoire n'en
contient aucun}, et il n'y en a jamais eu. Le mot y a deux emplois, tous deux légitimes et tous
deux sous mesure physique : la forme de Dirichlet d'une chaîne de Markov, et un brownien sans
dérive pour un temps de premier passage. Aucun prix, aucun actif répliqué, aucun changement de
mesure.
\vspace{1mm}

Un test \emph{est} conduit, mais il ne porte pas sur une martingale : c'est l'asymétrie du
corpus contre une loi nulle exacte, et son résultat est publié \textbf{et négatif}, $119$ contre
$128 \pm 27$, soit $z = -0{,}33$.
\begin{block}{\garder{} les deux paragraphes de précaution}
Ils coûtent quelques lignes et ferment définitivement une objection qu'un jury reposerait. Le
lien rhétorique qui suggérait un cadre commun entre les deux emplois est corrigé :
\textbf{ce qu'ils partagent est la mesure, pas la mathématique}.
\end{block}
\srcs{scripts~\texttt{40} et~\texttt{64}}
\end{frame}

\begin{frame}{C4. La sensibilité des paramètres}
\scriptsize
\textbf{Demandé.} Montrer la sensibilité aux paramètres, en particulier à la propagation.
\vspace{1mm}

\textbf{Fait, et la conclusion a dû être refaite avant d'être défendable.} Le tornado d'origine
comparait des \emph{longueurs de barre} sur des plages d'inégale vraisemblance, et son levier
« seuil » n'était pas un levier : le déplacer réajuste la loi de queue, donc il bougeait quatre
paramètres quand le gain en bougeait un. La normalisation qui rend les leviers comparables est
l'\textbf{élasticité}, même perturbation relative pour tous, un paramètre à la fois.
\begin{center}
\begin{tabular}{@{}l l r@{}}
\toprule
\textbf{Levier} & \textbf{Statut} & \textbf{Élasticité} \\
\midrule
Indice de queue $\xi$ & estimé & $3{,}79$ \\
Échelle $\sigma$ & estimé & $0{,}97$ \\
Fréquence $\lambda$ & estimé & $0{,}45$ \\
\textbf{Gain de propagation} $\boldsymbol{g}$ & \textbf{posé} & $\mathbf{0{,}37}$ \\
Seuil $u$ seul & choix de méthode & $0{,}03$ \\
\bottomrule
\end{tabular}
\end{center}
\vspace{-1mm}
\begin{block}{\garder, \emph{dans sa version renormalisée}}
\textbf{La conclusion tient et elle est plus forte qu'avant} : un rapport de dix entre la queue
et la propagation, qui ne dépend d'aucune convention d'affichage. L'ancienne formulation
comparait des amplitudes et pouvait être retournée en changeant une plage. Le seuil \emph{seul}
ne déplace que $+320$~M\euro{} quand l'indice de queue seul en déplace $-13\,993$ : tout ce
qu'on attribuait au seuil venait de la queue qu'il entraîne.
\end{block}
\srcs{scripts~\texttt{15}, \texttt{30} et~\texttt{76}}
\end{frame}

\begin{frame}{C5 et C6. La CTE, et le niveau de l'intervalle reporté}
\scriptsize
\textbf{C5, demandé.} Une CTE ou CVaR à côté de la VaR.
\vspace{1mm}

\textbf{La première réponse est réglementaire, et c'est par là qu'il faut commencer :} le régime
applicable définit le capital comme une VaR à $99{,}5\,\%$ sur un an, quand le Test suisse de
solvabilité retient une CTE à $99\,\%$. Se placer sous le premier régime est la raison
principale de ne pas substituer l'une à l'autre.
\vspace{1mm}

\textbf{Les mesures viennent ensuite, et elles renforcent.} CTE de $5\,734$, $12\,607$ et
$18\,183$ aux niveaux $95$, $99$ et $99{,}5\,\%$. \textbf{La CTE à $95\,\%$ est moins prudente
que la VaR réglementaire}, l'équivalence tombant à $\beta = 97{,}73\,\%$. Et sous la calibration
alternative $\hat\xi = 1{,}033 > 1$, l'espérance est infinie et \textbf{la mesure n'existe pas}.
\vspace{1mm}

\textbf{C6, demandé.} Reporter le SCR avec un intervalle à $95\,\%$.
\vspace{1mm}

\textbf{Fait sans rejouer aucun bootstrap.} L'indice de queue passerait à
$[0{,}2487\,;0{,}8765]$, le capital à $[2\,727\,;34\,942]$~M\euro.
\begin{block}{\garder{} la CTE, et \garder{} le niveau à $90\,\%$}
Une page pour un argument qui ferme la question : le meilleur rapport des deux calls. Pour le
niveau, la décision est prise et écrite : \textbf{rester à $90\,\%$}. La bascule ne déplace
aucune calibration et ne corrige pas la sous-couverture, qui reste de $3{,}2$ points aux deux
niveaux. \emph{Et ce déficit ne doit pas être dramatisé} : trois points sur un intervalle de
paramètre de forme à quatre-vingt-onze excès est l'ordre de grandeur attendu, la normalité
asymptotique y convergeant lentement. Il se publie, il ne disqualifie pas.
\end{block}
\srcs{scripts~\texttt{72} et~\texttt{73}}
\end{frame}

\begin{frame}{C7 et C8. La renormalisation par la taille, et les trois étages}
\scriptsize
\textbf{C7, demandé.} Renormaliser proportionnellement au SCR de marché.
\vspace{1mm}

\textbf{Testé, et la mesure conclut contre} : élasticité $0{,}204$ contre $1$ supposée. Le détail
est en slide H5, la demande étant la même vue des deux côtés.
\vspace{1mm}

\textbf{C8, demandé.} Chiffrer les trois étages d'incertitude, puis retirer le troisième.
\vspace{1mm}

\textbf{Chiffré, et le résultat va contre l'intuition de l'arbitrage} : l'étage de modèle est
\textbf{du même ordre} que celui de paramètre, facteur $9{,}21$ sur l'axe de la famille de queue
contre $8{,}93$, et tous deux \textbf{dominent largement} l'identification à $1{,}27$. Le
retirer rétrécirait donc l'affichage autant que retirer le paramètre, alors que l'incertitude ne
bougerait pas.
\vspace{0.5mm}

\emph{Formulation corrigée après relecture : dire \og le plus gros des trois \fg{} sur un écart
de $3\,\%$ entre $9{,}21$ et $8{,}93$ n'est pas soutenable à quatre graines. Ce qui est résolu
est le rapport aux $1{,}27$ de l'identification, pas l'ordre entre les deux premiers.}
\begin{block}{\garder{} le compromis, \compresser{} son exposé}
L'étage de modèle \textbf{sort de la bande reportée} mais devient un axe de sensibilité déclaré,
et reste publié deux fois. Le motif écrit est ce qui rend la décision défendable : les deux
premiers étages sont \emph{statistiques} donc portables par un intervalle, le troisième est
\emph{épistémique} et un choix de famille ne se moyenne pas. Ce motif tient en un paragraphe.
\end{block}
\srcs{scripts~\texttt{48}, \texttt{70} et~\texttt{71}}
\end{frame}

\section*{Le call avec Caroline : structure, et voies refermées}

\begin{frame}{C9. La conformité aux piliers 1 et 2 protège-t-elle le 3}
\scriptsize
\textbf{Demandé.} Être conforme aux piliers 1 et 2 rend indirectement un peu conforme au 3.
\vspace{1mm}

\textbf{Trois choses, dans cet ordre.}
\begin{itemize}\itemsep1pt
\item \textbf{Répondre que la dépendance manque serait faux} : elle existe déjà, de corrélation
$0{,}462$ entre deux piliers quelconques. Ce qui manque est sa \emph{structure}, échangeable au
lieu d'être propre à certains triplets ;
\item \textbf{l'omission ne coûte rien, et c'est mesuré} : en ajoutant un surcroît de
corrélation aux seules paires concernées, admissible jusqu'à $0{,}38$, la loi des configurations
bouge nettement mais le capital espéré ne bouge que de $2$~M\euro, soit $0{,}02\,\%$ ;
\item \textbf{la raison est structurelle, pas numérique} : le capital des $32$ configurations
s'ajuste par une forme \emph{additive} en indicateurs de pilier à $R^2 = 0{,}9945$, et
l'espérance d'une fonction additive ne dépend que des \textbf{marges}.
\end{itemize}
\begin{block}{\garder, \emph{avec la réserve désormais mesurée}}
\textbf{Le meilleur résultat des deux calls.} L'invariance n'est pas constatée sur un balayage,
elle est \emph{démontrée} : elle vaut pour \textbf{toute} structure à marges fixées, y compris
celles qu'on n'a pas essayées. Mais elle porte sur une \textbf{espérance}, et le capital est un
quantile : la slide suivante mesure ce que cela change.
\end{block}
\srcs{script~\texttt{69}}
\end{frame}

\begin{frame}{C9 \emph{(suite)}. L'objection qui manquait, et sa mesure}
\scriptsize
\textbf{L'argument d'additivité vaut pour une espérance, or le capital est un quantile.} Ce que
le calcul établissait est la \emph{moyenne des VaR conditionnelles} aux $32$ configurations. Un
lecteur qui demande le capital d'une entité dont l'état de conformité est incertain demande le
quantile de la loi \textbf{mélangée}, et la moyenne des quantiles n'est pas le quantile du
mélange. Les deux ont donc été calculés.
\vspace{1mm}

\begin{center}
\begin{tabular}{@{}l r r@{}}
\toprule
& \textbf{Niveau} & \textbf{Déplacement sur le balayage} \\
\midrule
Moyenne des VaR conditionnelles & $9\,739$~M\euro & $2$~M\euro \\
Quantile du mélange & $10\,643$~M\euro & $334$~M\euro \\
\midrule
\emph{bruit de tirage seul, 8 graines} & & \emph{étendue} $344$~M\euro \\
\bottomrule
\end{tabular}
\end{center}
\vspace{-1mm}
\begin{block}{\garder{} la distinction, et la formuler honnêtement}
\textbf{Deux enseignements.} Les deux objets diffèrent de $904$~M\euro, soit $9{,}3\,\%$ : le
mémoire doit dire lequel il publie, et il le dit désormais. Et sur le quantile,
\textbf{l'invariance n'est pas démontrée, elle est seulement non détectée} : le déplacement de
$334$~M\euro{} est du même ordre que le bruit de tirage, et il n'est pas monotone. La conclusion
de gestion ne bouge pas, ces $334$~M\euro{} valant $18\,\%$ de la bande d'identification ; c'est
le \emph{statut} de la conclusion qui passe de démontré à non détecté.
\end{block}
\srcs{scripts~\texttt{69} et~\texttt{77}}
\end{frame}

\begin{frame}{C10. Les défaillances simultanées et l'additivité des coûts}
\scriptsize
\textbf{Demandé.} Le cas où plusieurs piliers défaillent ensemble, et l'additivité de leurs
coûts.
\vspace{1mm}

\textbf{La prémisse est à corriger d'abord.} Ce n'est pas un cas laissé de côté, \textbf{c'est
la sortie du modèle}, et à l'état non conforme la défaillance multiple est \emph{majoritaire} :
$62{,}31\,\%$ des sinistres touchent plus d'un pilier, contre $31{,}15\,\%$ à l'état conforme.
\vspace{1mm}

\textbf{Ensuite, \og additivité \fg{} désigne trois choses à trois étages.} Hypothèse non testée
sur les coûts d'un sinistre ; quasi-additivité \emph{mesurée} en fonction des piliers ;
super-additivité \emph{mesurée} en fonction des canaux, $+35\,\%$. \textbf{Répondre \og le
modèle est super-additif \fg{} serait donc faux.}
\vspace{1mm}

\textbf{Enfin l'hypothèse est bornée} : un exposant de $0{,}70$ à $1{,}30$ déplace l'écart de
$9\,706$ à $20\,487$~M\euro, soit $15$ fois son bruit de $\pm 734$.
\begin{block}{\annexer{} le détail, \garder{} un paragraphe dans le corps}
Le corps doit garder la distinction des trois étages, qui évite une réponse fausse, et
l'asymétrie inattendue : élasticités $0{,}95$ à l'état non conforme contre $0{,}39$ à l'état
conforme. Le reste, la table complète des exposants et la loi du cardinal, vit très bien en
annexe.
\end{block}
\srcs{script~\texttt{74}}
\end{frame}

\begin{frame}{C11 et C12. Les séquences ordonnées, et les dates d'arrivée}
\scriptsize
\textbf{C11, demandé.} Exploiter les probabilités conditionnelles de séquence, du type
P1 $\rightarrow$ P2 $\rightarrow$ P3, pour reconstituer la direction.
\vspace{1mm}

\textbf{La voie a été suivie jusqu'au bout et se referme, pour trois raisons distinctes.} Les
séquences étaient \textbf{déjà dans le corpus} : composer les arêtes donne $9$ chemins et $4$
séquences distinctes. \textbf{La comparaison demandée a un support vide}, P2 n'émettant jamais,
$0$ fois sur $22$. \textbf{Et le test n'a aucune puissance}, pas faute d'échantillon : il exige
d'un même pilier qu'il soit atteint \emph{et} qu'il ait deux successeurs, et les deux manques
sont exclusifs dans ce corpus.
\vspace{1mm}

\textbf{C12, demandé.} Le processus des dates d'arrivée, et l'effet du regroupement.
\vspace{1mm}

\textbf{Non traité, et la raison est dans la donnée} : le coût d'un sinistre \emph{daté} n'est
jamais observé, la chronologie portant les dates et la conversion portant les coûts.
\begin{block}{\compresser{} le corps, \garder{} l'annexe}
Un résultat négatif se défend quand il livre un \textbf{critère de réouverture} précis, et
celui-ci le fait : un pilier atteint avec deux successeurs, répété $158$ à $589$ fois. À
conserver absolument : la chaîne \textbf{n'est pas sans mémoire}, l'auto-évitement gonflant la
loi du successeur de $1{,}211$, donc les triplets ne sont pas redondants avec les paires.
\end{block}
\srcs{script~\texttt{75}}
\end{frame}

\begin{frame}{C13. La méthode de Cooke en annexe, et le questionnaire}
\scriptsize
\textbf{Demandé.} Mettre la méthode et le questionnaire en annexe, peaufiner le questionnaire.
\vspace{1mm}

\textbf{Fait, et l'annexe a changé de sens.} Elle ne présente pas un protocole à venir mais
\textbf{un protocole préparé et non exécuté}, avec le chiffrage de la raison : un panel
faiblement calibré déplacerait le capital de $777$~M\euro{} \emph{sans qu'aucun signal ne
l'indique}, là où l'approche par bornes reste invariante. Le coût de l'abandon est nommé : P1 et
P4 restent à égalité, P1 en tête dans $43{,}6\,\%$ des cas, pour un écart de regret maximal de
$1{,}5\,\%$.
\vspace{1mm}

\textbf{Le questionnaire a été repris, et la reprise a évité un vrai dégât.} Une graine était
\textbf{fausse} : elle annonçait un indice de queue de $0{,}90$, valeur posée du pire cas, quand
la calibration donne $0{,}5954$. Dans la méthode de Cooke une graine fausse n'ajoute pas du
bruit, elle \textbf{inverse} les poids. Et la convention de direction était \emph{inversée} par
rapport au mémoire : une saisie fidèle aurait transposé la matrice.
\begin{block}{\garder}
Trois pages d'annexe bien employées : c'est un actif devant un jury qui demanderait pourquoi un
modèle dont la direction n'est pas identifiable ne recourt pas au jugement d'expert. La décision
reste \textbf{réversible sans coût de développement}.
\end{block}
\srcs{scripts~\texttt{26}, \texttt{30}, \texttt{31} et~\texttt{47}}
\end{frame}

\section*{Ce qui vous revient}

% HARNAIS-HORS-SCRIPT: le nombre de fichiers touches par le renommage est un compte de code
% source, pas une grandeur du modele. Aucun script ne le produit et aucun ne le devrait.

\begin{frame}{Les quatre décisions, tranchées comme je les signerais}
\scriptsize
Ces quatre points vous reviennent, mais laisser un arbitrage ouvert n'est pas neutre non plus.
Voici la position que je défendrais devant un comité, avec le motif professionnel qui la porte.
Elle reste à ratifier.
\vspace{1.5mm}
\begin{block}{1. Publier l'écart à quatre canaux : $\Delta = 14\,139 \pm 734$~M\euro}
Le $-53\,\%$ garde le \textbf{gain de propagation au niveau non conforme} pendant qu'il fait
varier le score : il décrit donc une entité conforme sur le papier qui propagerait comme une
non conforme. \textbf{Ce n'est pas un état du monde cohérent}, c'est une dérivée partielle
présentée comme un scénario, et un scénario incohérent ne se reporte pas. L'écart à quatre
canaux, lui, oppose deux états complets, décompose son interaction et publie son bruit.
\emph{Le $-53\,\%$ reste, réétiqueté en sensibilité au seul score de conformité.}
\end{block}
\vspace{-1.5mm}
\begin{block}{2. Rester à $90\,\%$, et publier la couverture atteinte}
La couverture atteinte manque le nominal de $-3{,}2$ points \textbf{aux deux niveaux} : passer à
$95\,\%$ ne corrige pas ce déficit, il le rebaptise. \textbf{On n'annonce pas un niveau de confiance qu'on n'atteint
pas}, et élargir la bande sans la rendre plus vraie donne une fausse impression de rigueur
accrue. Le correctif propre serait d'élargir la demi-largeur d'environ $10\,\%$, mais c'est une
recalibration que le gel interdit. \emph{Donc : $90\,\%$ annoncé, $86{,}8 \pm 0{,}8\,\%$ mesuré
et écrit à côté.} \textbf{Appliqué} : l'annexe portait encore \og le niveau reporté est
$95\,\%$ \fg{} alors que le corps imprime partout les bornes à $90\,\%$, et l'incohérence est
levée.
\end{block}
\srcs{scripts~\texttt{25}, \texttt{43}, \texttt{47}, \texttt{72} et~\texttt{74}}
\end{frame}

\begin{frame}{Les quatre décisions, tranchées \emph{(suite)}}
\scriptsize
\begin{block}{3. Comprimer l'illustration, garder la borne de validité}
Principe de proportionnalité : \textbf{la place accordée à un résultat doit suivre sa
fiabilité}. Le niveau d'entité est déclaré \emph{borne supérieure d'ordre de grandeur}, sous la
borne inférieure de validité de la descente d'échelle. Lui donner quatre pages d'illustration
invite le lecteur à s'en servir comme d'une mesure. \emph{Je comprime la partie illustrative à
une page et demie, et je garde intégralement les deux pages qui établissent la borne inférieure
de validité sur bilans réels : c'est une limite, et une limite ne se comprime pas.}
\end{block}
\vspace{-1.5mm}
\begin{block}{4. Ne pas renommer, mais verrouiller par un contrôle automatique}
Une collision de noms dans un moteur de calcul est un \textbf{risque de modèle}, pas une
question de style. Mais renommer vingt-sept fichiers sur un pipeline gelé, à quelques semaines
du rendu, introduit un risque d'erreur supérieur au risque qu'on veut traiter. \textbf{Quand la
remédiation coûte plus que le risque, on ne remédie pas, on contrôle.} \emph{Donc : pas de
renommage avant la soutenance, et une assertion à l'import qui vérifie que les deux constantes
valent ce qu'elles doivent valoir. Elle ne déplace aucun nombre et rend la confusion
impossible.}
\end{block}
\vspace{-1.5mm}
\begin{block}{Le fil commun des quatre}
Aucune ne consiste à corriger le modèle. Trois consistent à \textbf{déclarer plutôt qu'à
rebaptiser}, et la quatrième à \textbf{contrôler plutôt qu'à remédier}. C'est la posture qu'un
gel de calibration impose, et c'est celle qui se défend devant un jury.
\end{block}
\end{frame}

\begin{frame}{Ce que je propose de garder, et ce que je propose de réduire}
\scriptsize
\begin{block}{À garder sans discussion}
La table complète des leviers et ses trois lectures. \textbf{L'allocation aux cinq piliers avec
ses deux réserves}, qui tient en une page. Le sens de la borne du retour. La CTE et son argument
d'existence. La sensibilité qui montre que la propagation est le levier le moins sensible.
L'invariance à la structure de dépendance, qui est \emph{démontrée} et non constatée. Le
protocole d'élicitation préparé et non exécuté.
\end{block}
\vspace{-1.5mm}
\begin{block}{À réduire, sans rien perdre d'essentiel}
Les scénarios par taille, de six pages à une et demie. L'exposé des trois étages, dont seul le
motif épistémique doit rester. Les défaillances simultanées, dont le détail vit en annexe. Les
séquences ordonnées, dont le corps peut se limiter au support vide et au critère de réouverture.
\end{block}
\vspace{-1.5mm}
\begin{block}{La règle qui gouverne tout le reste}
\textbf{Aucune limite déclarée n'est proposée au retrait.} Ce mémoire tient parce qu'il publie
ses propres réserves : la non-transitivité réfutée par son auteur, le chiffre central requalifié
en borne supérieure, la borne inférieure de validité publiée. Gagner une page en rendant une
réserve moins visible serait le plus mauvais échange possible.
\end{block}
\end{frame}

\begin{frame}{Une précaution sur le contrôle lui-même}
\scriptsize
Chaque chiffre de ce document est vérifié automatiquement : un harnais cherche chaque nombre
publié dans la sortie des scripts que sa section cite. C'est ce qui permet d'affirmer qu'aucun
nombre n'est calculé à la main. Mais il faut dire ce que ce contrôle \emph{ne} fait pas, sans
quoi le taux qu'il produit serait surinterprété.
\vspace{1.5mm}
\begin{block}{Le harnais mesure une couverture, pas une exactitude}
Il vérifie qu'un nombre publié \textbf{existe} dans la sortie du script cité. Il ne vérifie ni
qu'il est au \textbf{bon endroit}, ni qu'il veut dire ce que la \textbf{phrase prétend}. Un
nombre confirmé peut donc être mal commenté, et un taux élevé ne signifie pas que le mémoire est
juste dans la même proportion.
\end{block}
\vspace{-1.5mm}
\begin{block}{Ce qu'il attrape réellement, et c'est déjà beaucoup}
Les chiffres \textbf{périmés}, ceux qu'aucun calcul actuel ne reproduit, et les grandeurs citées
sans être imprimées par aucun script. Sur les deux calls, c'est ainsi qu'ont été trouvés une
borne d'intervalle qui n'appartenait à aucune méthode, une graine d'étalonnage fausse dans le
questionnaire d'expert, et deux nombres de mise en page comptés comme des résultats.
\end{block}
\vspace{-1.5mm}
\begin{block}{La conséquence pratique}
Ce que garantit le dispositif est la \textbf{traçabilité}, pas la justesse d'interprétation.
Cette dernière ne s'obtient que par la relecture, et c'est précisément ce que ces deux calls
apportent.
\end{block}
\end{frame}

\end{document}
